\anexo{Encuesta: instrumento y resultados agregados}\label{annex:user-research}

Este anexo reproduce el cuestionario, las frecuencias agregadas de las preguntas cerradas y las respuestas a la pregunta abierta en que se apoya la Subsubsección~\ref{subsec:encuesta}. Las preguntas se transcriben tal como figuraban en el relevamiento, con sus redundancias incluidas.

\section*{Encuesta: formulario completo y respuestas}

La encuesta, titulada \textit{Encuesta sobre análisis y gestión de inversiones}, se administra con Google Forms y queda abierta del 24 de mayo al 14 de agosto de 2026; en ese período se registran 190 respuestas completas. A continuación se transcribe el formulario completo tal como lo vieron los encuestados. La Tabla~\ref{tab:encuesta-distribucion} presenta la distribución de respuestas de cada pregunta cerrada, la Tabla~\ref{tab:encuesta-temas} la codificación temática de la pregunta abierta y la Tabla~\ref{tab:encuesta-abierta} la transcripción de sus respuestas.

\subsection*{Formulario completo}

El formulario incluye opciones únicas, selecciones múltiples, escalas y una pregunta abierta.

\begin{enumerate}
    \item \textbf{Rango de edad.} Opción única: 18--24; 25--34; 35--44; 45+.
    \item \textbf{¿Cuál de estas situaciones te representa mejor?} Opción única: Invierto actualmente de forma activa; Invierto, pero hago seguimiento ocasional; No invierto actualmente, pero me interesa empezar; No invierto y no me interesa por ahora.
    \item \textbf{Tu nivel de conocimiento financiero.} Escala de 1 (muy bajo) a 5 (muy alto).
    \item \textbf{¿Qué tipos de activos tenés actualmente o te interesaría tener?} Selección múltiple: Acciones / CEDEARs; ETFs; Criptomonedas; Stablecoins / DeFi; Bonos / FCI / Plazo fijo; Dólar (MEP, CCL, físico); Otros / Ninguno todavía.
    \item \textbf{¿En cuántas plataformas distintas (brokers + exchanges) tenés activos o seguís inversiones?} Opción única: Ninguna todavía; 1; 2; 3; 4 o más.
    \item \textbf{¿Cómo hacés hoy el seguimiento de tu portafolio?} Opción única: Uso la app de cada plataforma por separado; Llevo un Excel / Google Sheets propio; Uso una app de tracking o portfolio management; Lo sigo de forma poco ordenada; No sigo un portafolio todavía.
    \item \textbf{¿Sentís que la información que necesitás para analizar tus inversiones está repartida entre muchas fuentes?} Escala de 1 (para nada) a 5 (totalmente).
    \item \textbf{¿Cuáles son tus principales dificultades al gestionar o analizar tu portafolio? (hasta 3)} Selección múltiple: Tener visión consolidada de todo; Saber si estoy bien diversificado; Entender mi exposición real (sector, región, moneda); Interpretar contexto macro y geopolítico; Seguir el mundo cripto / DeFi; Saber cuándo rebalancear; Falta de tiempo para investigar; No sé qué fuentes son confiables; Traducir mucha información en conclusiones claras.
    \item \textbf{¿Alguna vez te enteraste tarde de algo (noticia, evento) que habría cambiado una decisión de inversión?} Opción única: Sí, varias veces; Sí, alguna vez; No estoy seguro; No.
    \item \textbf{¿Conocés o usaste plataformas de prediction markets (por ejemplo, Polymarket)?} Opción única: Sí, las uso activamente; Sí, las conozco pero no las uso; No las conocía.
    \item \textbf{¿Te interesaría que una herramienta de análisis de portafolio incorpore señales basadas en probabilidades de eventos futuros (por ejemplo, tasas, inflación, elecciones o regulaciones)?} Escala de 1 (nada interesado) a 5 (muy interesado).
    \item \textbf{Si existiera una plataforma web que te permita cargar tu portafolio (acciones + ETFs + cripto) y te dé análisis integrado con contexto geopolítico y señales de mercado, ¿qué tan interesado estarías?} Escala de 1 (nada interesado) a 5 (muy interesado).
    \item \textbf{De estas funcionalidades, marcá las 3 más valiosas para vos.} Selección múltiple: Vista consolidada de todos mis activos; Análisis de diversificación y exposición real; Análisis fundamental básico de mis activos; Alertas sobre eventos geopolíticos relevantes para mi portafolio; Señales del ecosistema cripto / DeFi / prediction markets; Recomendaciones explicables (que muestren el por qué); Sugerencias de rebalanceo.
    \item \textbf{¿Qué tan importante te parece que una herramienta de IA explique por qué te muestra una observación o recomendación?} Escala de 1 (nada importante) a 5 (muy importante).
    \item \textbf{¿Usás o usaste alguna herramienta similar?} Opción única: Sí, la uso actualmente; La usé y la dejé; Las conozco pero no las uso; No conocía ninguna.
    \item \textbf{¿Qué sentís que le falta a las herramientas que usás hoy para analizar mejor tu portafolio?} Pregunta abierta.
    \item \textbf{¿Pagarías por una herramienta así?} Opción única: Sí, mensualmente; Solo si tiene versión gratuita con funciones pagas; Solo si es 100\,\% gratis; No la usaría.
\end{enumerate}

\subsection*{Distribución de respuestas}

Los porcentajes toman N=190 como denominador; en las preguntas de selección múltiple pueden sumar más de 100\,\%, y las pequeñas diferencias en las sumas se deben al redondeo. Fuente: elaboración propia sobre las frecuencias documentadas en el relevamiento.

{
\small
\singlespacing
\setlength{\tabcolsep}{5pt}
\renewcommand{\arraystretch}{1.15}
\begin{longtable}{p{9.3cm} r r}
	\caption{Distribución de respuestas (N=190). Fuente: elaboración propia.}
	\label{tab:encuesta-distribucion}\\
	\hline
	\textbf{Pregunta / Opción} & \textbf{n} & \textbf{\%} \\
	\hline
	\endfirsthead
	\hline
	\textbf{Pregunta / Opción} & \textbf{n} & \textbf{\%} \\
	\hline
	\endhead
	\hline
	\endfoot
	\multicolumn{3}{l}{\textbf{1. Rango de edad}} \\
	\quad 18--24 & 58 & 30,5 \\
	\quad 25--34 & 77 & 40,5 \\
	\quad 35--44 & 37 & 19,5 \\
	\quad 45+ & 18 & 9,5 \\
	\multicolumn{3}{l}{\textbf{2. Situación frente a la inversión}} \\
	\quad Invierto actualmente de forma activa & 85 & 44,7 \\
	\quad Invierto, pero hago seguimiento ocasional & 57 & 30,0 \\
	\quad No invierto actualmente, pero me interesa empezar & 36 & 18,9 \\
	\quad No invierto y no me interesa por ahora & 12 & 6,3 \\
	\multicolumn{3}{l}{\textbf{3. Nivel de conocimiento financiero (media: 3,17)}} \\
	\quad 1 & 17 & 8,9 \\
	\quad 2 & 33 & 17,4 \\
	\quad 3 & 60 & 31,6 \\
	\quad 4 & 61 & 32,1 \\
	\quad 5 & 19 & 10,0 \\
	\multicolumn{3}{l}{\textbf{4. Tipos de activos que tiene o le interesaría tener (múltiple)}} \\
	\quad Acciones / CEDEARs & 129 & 67,9 \\
	\quad ETFs & 105 & 55,3 \\
	\quad Dólar (MEP, CCL, físico) & 103 & 54,2 \\
	\quad Bonos / FCI / Plazo fijo & 90 & 47,4 \\
	\quad Criptomonedas & 90 & 47,4 \\
	\quad Stablecoins / DeFi & 49 & 25,8 \\
	\quad Otros / Ninguno todavía & 13 & 6,8 \\
	\multicolumn{3}{l}{\textbf{5. Cantidad de plataformas (brokers + exchanges)}} \\
	\quad Ninguna todavía & 24 & 12,6 \\
	\quad 1 & 24 & 12,6 \\
	\quad 2 & 44 & 23,2 \\
	\quad 3 & 37 & 19,5 \\
	\quad 4 o más & 61 & 32,1 \\
	\multicolumn{3}{l}{\textbf{6. Forma de seguimiento del portafolio}} \\
	\quad Uso la app de cada plataforma por separado & 58 & 30,5 \\
	\quad Llevo un Excel / Google Sheets propio & 41 & 21,6 \\
	\quad No sigo un portafolio todavía & 41 & 21,6 \\
	\quad Uso una app de tracking o portfolio management & 28 & 14,7 \\
	\quad Lo sigo de forma poco ordenada & 22 & 11,6 \\
	\multicolumn{3}{l}{\textbf{7. Percepción de información repartida entre muchas fuentes (media: 3,51)}} \\
	\quad 1 & 13 & 6,8 \\
	\quad 2 & 30 & 15,8 \\
	\quad 3 & 38 & 20,0 \\
	\quad 4 & 66 & 34,7 \\
	\quad 5 & 43 & 22,6 \\
	\multicolumn{3}{l}{\textbf{8. Principales dificultades (múltiple, hasta 3)}} \\
	\quad Tener visión consolidada de todo & 75 & 39,5 \\
	\quad Falta de tiempo para investigar & 68 & 35,8 \\
	\quad Saber si estoy bien diversificado & 60 & 31,6 \\
	\quad Seguir el mundo cripto / DeFi & 52 & 27,4 \\
	\quad Entender mi exposición real (sector, región, moneda) & 51 & 26,8 \\
	\quad Traducir mucha información en conclusiones claras & 50 & 26,3 \\
	\quad Interpretar contexto macro y geopolítico & 49 & 25,8 \\
	\quad No sé qué fuentes son confiables & 45 & 23,7 \\
	\quad Saber cuándo rebalancear & 38 & 20,0 \\
	\multicolumn{3}{l}{\textbf{9. Se enteró tarde de una noticia o evento relevante}} \\
	\quad Sí, alguna vez & 75 & 39,5 \\
	\quad Sí, varias veces & 56 & 29,5 \\
	\quad No estoy seguro & 36 & 18,9 \\
	\quad No & 23 & 12,1 \\
	\multicolumn{3}{l}{\textbf{10. Conocimiento de prediction markets}} \\
	\quad No las conocía & 103 & 54,2 \\
	\quad Sí, las conozco pero no las uso & 71 & 37,4 \\
	\quad Sí, las uso activamente & 16 & 8,4 \\
	\multicolumn{3}{l}{\textbf{11. Interés en señales basadas en probabilidades de eventos (media: 3,53)}} \\
	\quad 1 & 11 & 5,8 \\
	\quad 2 & 30 & 15,8 \\
	\quad 3 & 45 & 23,7 \\
	\quad 4 & 55 & 28,9 \\
	\quad 5 & 49 & 25,8 \\
	\multicolumn{3}{l}{\textbf{12. Interés en la plataforma propuesta (media: 3,92)}} \\
	\quad 1 & 7 & 3,7 \\
	\quad 2 & 13 & 6,8 \\
	\quad 3 & 32 & 16,8 \\
	\quad 4 & 74 & 38,9 \\
	\quad 5 & 64 & 33,7 \\
	\multicolumn{3}{l}{\textbf{13. Funcionalidades más valiosas (múltiple, hasta 3)}} \\
	\quad Vista consolidada de todos mis activos & 116 & 61,1 \\
	\quad Análisis de diversificación y exposición real & 89 & 46,8 \\
	\quad Recomendaciones explicables (que muestren el por qué) & 88 & 46,3 \\
	\quad Alertas sobre eventos geopolíticos relevantes & 78 & 41,1 \\
	\quad Sugerencias de rebalanceo & 72 & 37,9 \\
	\quad Análisis fundamental básico de mis activos & 64 & 33,7 \\
	\quad Señales del ecosistema cripto / DeFi / prediction markets & 62 & 32,6 \\
	\multicolumn{3}{l}{\textbf{14. Importancia de la explicabilidad en herramientas de IA (media: 3,64)}} \\
	\quad 1 & 9 & 4,7 \\
	\quad 2 & 26 & 13,7 \\
	\quad 3 & 41 & 21,6 \\
	\quad 4 & 62 & 32,6 \\
	\quad 5 & 52 & 27,4 \\
	\multicolumn{3}{l}{\textbf{15. Uso de herramientas similares}} \\
	\quad No conocía ninguna & 81 & 42,6 \\
	\quad Las conozco pero no las uso & 43 & 22,6 \\
	\quad La usé y la dejé & 42 & 22,1 \\
	\quad Sí, la uso actualmente & 24 & 12,6 \\
	\multicolumn{3}{l}{\textbf{17. Disposición a pagar}} \\
	\quad Solo si tiene versión gratuita con funciones pagas & 95 & 50,0 \\
	\quad Sí, mensualmente & 49 & 25,8 \\
	\quad Solo si es 100\,\% gratis & 30 & 15,8 \\
	\quad No la usaría & 16 & 8,4 \\
\end{longtable}
}

\subsection*{Respuestas a la pregunta abierta}

La pregunta 16 recibió 190 respuestas, de las cuales 27 no aportan contenido: un guion, un punto o expresiones como \enquote{no sé} y \enquote{no uso}. Las 163 restantes se codificaron por tema con el programa \texttt{scripts/codificar\_respuestas\_abiertas.py} del repositorio del Proyecto, que asigna cada tema según las expresiones presentes en el texto; la asignación se revisó respuesta por respuesta, y una respuesta puede corresponder a más de un tema. La Tabla~\ref{tab:encuesta-temas} presenta las frecuencias, con porcentajes sobre las 163 respuestas sustantivas. La Tabla~\ref{tab:encuesta-abierta} transcribe las 190 respuestas en el orden en que se registraron, sin corregir ortografía ni redacción; algunos textos se repiten de forma idéntica o con variaciones mínimas.

\begin{table}[ht]
\centering\small\renewcommand{\arraystretch}{1.15}
\caption{Temas de las respuestas a la pregunta abierta (163 respuestas sustantivas). Fuente: elaboración propia.}\label{tab:encuesta-temas}
\begin{tabular}{>{\raggedright\arraybackslash}p{10.2cm}rr}\hline
\textbf{Tema} & \textbf{n} & \textbf{\%}\\\hline
Contexto argentino: inflación, tipo de cambio e instrumentos locales & 46 & 28,2 \\
Consolidación e integración entre plataformas & 39 & 23,9 \\
Noticias, alertas y efecto del contexto sobre la cartera & 34 & 20,9 \\
Carga y actualización automáticas & 23 & 14,1 \\
Lenguaje accesible y facilidad de uso & 18 & 11,0 \\
Exposición real, diversificación y composición de fondos & 18 & 11,0 \\
Criptoactivos, DeFi y mercados de predicción & 18 & 11,0 \\
Explicación y fundamento de las observaciones & 14 & 8,6 \\
Falta de tiempo para seguir y procesar la información & 10 & 6,1 \\
Otros temas & 6 & 3,7 \\
\hline

\end{tabular}
\end{table}
\FloatBarrier

{
\footnotesize
\singlespacing
\setlength{\tabcolsep}{5pt}
\renewcommand{\arraystretch}{1.1}
\begin{longtable}{r >{\raggedright\arraybackslash}p{13.6cm}}
	\caption{Respuestas a la pregunta abierta, en orden de registro (N=190). Fuente: elaboración propia.}
	\label{tab:encuesta-abierta}\\
	\hline
	\textbf{N.º} & \textbf{Respuesta} \\
	\hline
	\endfirsthead
	\hline
	\textbf{N.º} & \textbf{Respuesta} \\
	\hline
	\endhead
	\hline
	\endfoot
	1 & - \\
2 & No lo sé \\
3 & No uso \\
4 & No conozco ninguna \\
5 & Que se agreguen automáticamente las operaciones que hago en cualquiera de mis cuentas. No tener que cargarlas manualmente. Por otro lado, jamás confiaría en las recomendaciones de la IA, deben tener sustento en algún reserach que ya exista, incluir fuentes, views de equipos de research de distintos brokers puede ser interesante \\
6 & Conocer la ganancia real. En Argentina muchas de las inversiones son en pesos y debido a la inflación es complicado saber si vas ganando o perdiendo. Con los CEDEARs también es un problema ya que su valor está ajustado con el dólar CCL dificultando calcular la ganancia de la inversión \\
7 & Aparte de los analyst ratings, información de porque creen que el activo subirá o bajará. Si tengo una IA que me ayude con esto. Sería increíble \\
8 & Mejor analisis de que compre y que vendi para saber la ganancia real de mi portafolio \\
9 & Análisis de los componentes de cada ETF/MF para conocer exposición real \\
10 & Contexto inflacionario y geopolitico \\
11 & No uso ninguna herramienta pero el proyecto me parece una buena idea. \\
12 & personalización y el conocimiento profundo del negocio \\
13 & Poder tener todo en una sola plataforma \\
14 & Hacerse mas visibles. \\
15 & Resumen de noticias relacionadas a mis inversiones. \\
16 & Recomendaciones explicables \\
17 & Integración con otras plataformas \\
18 & No lo sé \\
19 & - \\
20 & Mas claridad y fácil utilización \\
21 & No conozco mucho de inversión \\
22 & Analisis \\
23 & Compra venta mas facil, mejor ux \\
24 & - \\
25 & confidencialidad \\
26 & Que puedan comparar exposición real \\
27 & No conozco ninguna \\
28 & No aplica \\
29 & Ni idea \\
30 & Tenes todo consolidado en un lugar y a partir de ahí tomar decisiones. \\
31 & . \\
32 & \_ \\
33 & Me cuesta tener todo junto en un solo lugar. Uso varias apps y pierdo tiempo consolidando datos manualmente cada semana. \\
34 & facilidad de conectar con mi portafolio actual \\
35 & No uso ese tipo de herramientas \\
36 & el problema es que tengo todo en pesos y no se bien como estoy expuesto al dolar en total. \\
37 & no encuentro nada que ande bien para el mercado argentino, todo esta pensado para afuera \\
38 & tiempo. no tengo tiempo de analizar todo yo solo. \\
39 & no encuentro nada que ande bien para el mercado argentino, todo esta pensado para afuera \\
40 & tmb me pasa que me entero de noticias importantes cuando ya es tarde lamentablemente \\
41 & cada vez que pasa algo macro tardo horas en entender si me afecta a mi o no. \\
42 & osea tengo cripto en binance y cedears en balanz y no hay forma de unirlo todo \\
43 & sigo tips de un grupo de telegram. se que no es lo ideal pero no encontre algo mejor. \\
44 & osea tengo cripto en binance y cedears en balanz y no hay forma de unirlo todo \\
45 & Necesito poder ver mi exposicion real por moneda, sector y region en un solo lugar. Hoy eso requiere mucho trabajo manual. \\
46 & el problema es que tengo todo en pesos y no se bien como estoy expuesto al dolar en total. \\
47 & las herraminetas que probe son muy tecnicas, necesito algo mas acccesible para mi nivel \\
48 & perdí plata por no enterarme a tiempo de una noticia, me hubiera venido bien una alerta \\
49 & que explique las cosas en criollo y no en terminos tecnicos. \\
50 & me falta una vista que me diga: esto es lo que tenes, esto es lo que rinde, esto es lo que riesgo. \\
51 & Lo que mas valoro es la explicabilidad. No alcanza con que me digan que hacer, necesito entender el razonamiento detras. \\
52 & Uso varias billeteras y plataformas y es un caos cuando quiero saer cuánto tengo \\
53 & Integración con IOL y Balanz como mínimo. \\
54 & No aplica \\
55 & tengo todo en DeFi y no hay ninguna app que integre eso con mis CEDEARs en Balanz. \\
56 & Mi problema es que invierto pero no entiendo bien lo que hago necesito más ayuda \\
57 & Análisis de la exposición real considerando los subyacentes de los CEDEARs y ETFs, no solo el valor nominal de las posiciones. \\
58 & Tengo cuenta en Balanz en IOL y en Lemon y cada uno me muestra sus números pero nunca veo el cuadro completo \\
59 & Me pasa que tengo FCI, CEDEARs y dólar MEP en distintos lados y nunca veo el total. \\
60 & Me falta una herramienta que integre el análisis fundamental con el contexto macroeconómico de forma automática y comprensible. \\
61 & Lo que más falta en el mercado argentino es una herramienta que trate las CEDEARs como lo que son: exposición a acciones extranjeras con riesgo cambiario local. \\
62 & sin comentarios \\
63 & Me cuesta saber si estoy bien diversificado o si en realidad estoy súper concentrado en algo. \\
64 & Alertas de noticias relevantes para mis posiciones. \\
65 & Automatizar la carga de operaciones. \\
66 & Que no sea tan técnico. \\
67 & No sé \\
68 & Visión consolidada de todo, eso es lo que falta. \\
69 & El análisis de liquidez de mis posiciones es algo que nunca contemplo bien. No sé cuánto tiempo me llevaría salir de cada activo si necesitara el dinero. \\
70 & Lo de la ganancia real en Argentina es muy complicado por la inflación y el tipo de cambio. \\
71 & Estaría bueno que me sugiera qué hacer cuando algo en el mercado cambia mucho. \\
72 & no entiendo nada de macro pero igual invierto, necesito que alguien me explique \\
73 & Uso varias billeteras y plataformas y es un caos cuando quiero saer cuánto tengo \\
74 & Lo que falta es algo que integre on-chain con off-chain en un solo dashboard. \\
75 & no se me ocurre nada \\
76 & El tema del tipo de cambio en Argentina complica todo. Uno no sabe bien si gano o no \\
77 & Alertas en tiempo real. \\
78 & no uso herramientas \\
79 & Me gustaría entender mejor los ETFs que tengo, qué hay adentro de cada uno. \\
80 & Me fala contexto macro. Sé que sube o baja pero no entiendo por q impacta en mí \\
81 & Me falta una vista que me diga: esto es lo que tenés, esto es lo que rinde, esto es lo que arriesgás. \\
82 & El seguimiento del ecosistema DeFi en conjunto con activos tradicionales es algo que ninguna herramienta local ofrece de forma integrada. \\
83 & Re complicado seguir todo junto. Necesito algo que lo haga automático ya \\
84 & Lo que más me molesta es que cuando quiero saber cuánto tengo en total tengo que abrir todo y calcularlo a mano \\
85 & Uso excel pero no tengo precios en tiempo real. \\
86 & Que me diga cuánto tengo en dólares reales, no nominales. \\
87 & Lo que más me falta es contexto. No solo precios sino por qué están subiendo o bajando. \\
88 & Mejor UX en general. \\
89 & Capaz me falta una herramienta que lo explique más simple, soy re principiante todavía \\
90 & Que me diga cuánto gané de verdad descontando inflación, eso es lo q más necesito \\
91 & Cada vez que suben las tasas no entiendo si eso me conviene o no, es medio confuso todo \\
92 & Tiempo y claridad. Eso es lo que más me falta para tomar mejores decisiones de inversión. \\
93 & Lo que más me molesta es no tener alertas cuando hay noticias que me pueden afectar. \\
94 & Lo que busco es algo que me traduzca la macro global en impacto concreto en mi portafolio. \\
95 & Igual lo sigo bastante poco la verdad, me falta tiempo y ganas de organizar todo \\
96 & Nada que sirva para el contexto argentino. \\
97 & Uso un excel que hice yo pero es un quilombo actualizarlo. Necesito algo automático. \\
98 & Todo muy fragmentado, abro 4 apps distintas y nunca veo el total real \\
99 & Información contextualizada para Argentina. \\
100 & N/A \\
101 & El excel me quedó chico pero no encontré nada mejor para el mercado local \\
102 & Probé apps pero ninguna entiende el mercado local. \\
103 & Integración con IOL y Balanz como mínimo. \\
104 & No entiendo bien mi exposición real, creo q tengo más riesgo del que pienso \\
105 & Poder ver todo en un solo lugar. \\
106 & Tengo un excel con 8 pestañas y igual no me alcanza. \\
107 & Tmbn me pasa que veo una noticia y no sé si me afecta a mí o no, es medio frustrante \\
108 & no lo sé \\
109 & ninguna \\
110 & Perdí plata por no enterarme a tiempo de una noticia, me hubiera venido bien una alerta \\
111 & Las herramientas existentes no contemplan el contexto argentino: inflación, tipo de cambio, riesgo país. Eso hace que el análisis quede incompleto. \\
112 & Me pasa seguido que veo algo en Twitter y no sé si es relevante para mis posiciones o es solo ruido y termino no haciendo nada \\
113 & Análisis de riesgo sencillo de entender. \\
114 & Me gustaría saber cuándo una noticia geopolítica me impacta directamente. \\
115 & Quería empezar a invertir pero no sé ni por dónde, todas las apps son muy confusas \\
116 & nose \\
117 & El análisis de ETFs debería incluir la composición interna para conocer la exposición real por sector y región, no solo el precio del ETF. \\
118 & Que las operaciones se carguen solas, cargarlas a mano es lo que más abandono. \\
119 & El problema es que la info está en todos lados y no tengo tiempo de procesarla toda y sacar una conclusión clara \\
120 & Lo que más valoro es la explicabilidad. No alcanza con que me digan qué hacer, necesito entender el razonamiento detrás. \\
121 & Lo de las CEDEARs es re complicado porque no sé si ganás en dólares de verdad o es todo dólar CCL \\
122 & La visualización de la exposición por moneda es algo que ninguna herramienta local ofrece de forma clara y automática. \\
123 & La idea de tener todo integrado con alertas me parece lo más útil de todo lo que mencionan. \\
124 & el problema es que tengo todo en pesos y no sé bien cómo estoy expuesto al dólar en total. \\
125 & Me cuesta mucho entender cuánto de mi portafolio está dolarizado de verdad. \\
126 & Cuando veo que el dólar se mueve mucho no sé si mis inversiones me cubren o estoy perdiendo plata real \\
127 & nada \\
128 & Tengo un FCI en Galicia y CEDEARs en IOL y dólares en Brubank y es imposible ver todo junto sin abrir tres apps \\
129 & No entiendo por qué ninguna app integra CEDEARs con cripto con FCI es como si los brokers no quisieran que vieras todo junto \\
130 & Me gustaría poder ver mi ganancia ajustada por inflación, no solo en pesos nominales. \\
131 & Ni idea \\
132 & El contexto argentino es muy específico y las apps gringas no lo contemplan para nada. \\
133 & No encuntro nada que funcione para el mercado local, todo está pensado para afuera \\
134 & Que integre cripto con lo tradicional. Eso básicamente. \\
135 & Usé Personal Portfolio en su momento pero no contemplaba el mercado argentino y lo dejé. \\
136 & El principal problema es la fragmentación. Tengo activos en distintas plataformas y no hay forma de ver el total consolidado sin hacerlo manualmente. \\
137 & Que muestre rendimientos ajustados por inflación y riesgo país. \\
138 & Se me va todo el finde actualizando el excel, re pajero eso \\
139 & Tmb me pasa que me entero de noticias importantes cuando ya es tarde lamentablemente \\
140 & Cuando sale una noticia macro siempre me pregunto si me afecta, y me lleva horas entenderlo. \\
141 & Cuando ocurren eventos geopolíticos relevantes no tengo forma de saber rápidamente cuáles de mis posiciones pueden verse afectadas. \\
142 & Me interesaría poder comparar el rendimiento de mi portafolio contra benchmarks relevantes para el mercado argentino, no solo el S\&P. \\
143 & Las herraminetas que probé son muy técnicas, necesito algo más accesible para mi nivel \\
144 & Qué molesto tener que ir app por app para ver cómo estoy. Qué paja \\
145 & La verdad no tengo ni idea de cuándo rebalancear, lo hago bastante rapido \\
146 & Probé Bloomberg y otras apps y nada funciona bien para el mercado argentino hay que adaptar todo a mano \\
147 & Cargo todo a mano y siempre lo termino abandonando. \\
148 & Que tenga en cuenta la inflación argentina. \\
149 & Que expliquen qué hay adentro de cada CEDEAR, no solo el precio. \\
150 & Me cueta mucho encontrar info confiable, hay mucho influencer dando consejos sin saber \\
151 & El seguimiento de prediction markets no está integrado en ningún análisis de portafolio serio que yo conozca. \\
152 & Lo que necesito es algo que me diga en qué estoy invertido de verdad sector moneda riesgo todo junto \\
153 & Me cansé de los grupos de Telegram con tips, necesito algo más serio pero que se entienda \\
154 & Re útil sería una cosa que te avise cuando algo importante pasa y te diga si te afecta o no \\
155 & Lo de los prediction markets es interesante pero nadie lo integra con el portafolio tradicional. \\
156 & Recomendaciones con fundamento, no solo ''comprá esto''. \\
157 & Seguir el ecosistema cripto requiere como 5 apps distintas y sigue siendo incompleto. \\
158 & No encuentro nada que ande bien para el mercado argentino, todo está pensado para afuera \\
159 & Una cosa que valoro mucho es que me expliquen el razonamiento, no solo la recomendación. \\
160 & Que explique en criollo y no en términos técnicos. \\
161 & Me muevo mucho entre chains y es imposible seguir todo con una sola herramienta. \\
162 & Me gustaría integrar mi cuenta de Lemon con lo q tengo en IOL Investor Pro \\
163 & El yield farming lo trackeo a mano, no encontré nada automático. \\
164 & La falta de integración entre plataformas es el principal cuello de botella. Debería poder conectar mis cuentas de Balanz, IOL y Binance en un solo lugar. \\
165 & No sé cuánto de lo que gano es real y cuánto es inflación. \\
166 & Sigo un grupo de Telegram. No es lo ideal pero no encontré algo mejor. \\
167 & Igual sigo todo por Twitter más que por las apps, no sé si es lo mejor jaja \\
168 & Cada semana hago el mismo ritual de exportar datos de cada plataforma pegarlo en el excel y calcular el total es una pérdida de tiempo enorme \\
169 & Me gustaría una herramienta que contextualice cada evento macroeconómico en función de mi portafolio específico, no de forma genérica. \\
170 & Abrí cuenta en tres brokers y ahora no sé ni cuánto tengo en total \\
171 & Cada vez que leo noticias de economía argentina no sé si actuar o quedarme quieto porque no entiendo bien el impacto \\
172 & Necesito poder distinguir entre ganancia nominal y ganancia real. Hoy todo se mezcla. \\
173 & Probé Delta y Monarch y nada, ninguna soporta CEDEARs bien \\
174 & La mayoría de portfolio trackers no soportan tokens en Polygon o Arbitrum, es frustrante. \\
175 & Cada vez que pasa algo macro tardo horas en entender si me afecta a mí o no. \\
176 & Me pasa que leo cosas contradictorias en distintos medios y no sé a quién creerle para tomar decisiones \\
177 & Ojalá sincronización con Mercado Pago y Belo tmbn, ahí tengo plata parada \\
178 & Lo que más me cuesta es saber cuándo y cómo rebalancear. Nadie te lo explica bien. \\
179 & Me custa consolidar todo. Uso tres apps y ninguna se conecta con la otra \\
180 & El mayor problema es k todo está separado y no hay forma de verlo junto fácilmente \\
181 & Las recomendaciones que encuentro en internet son genéricas. Necesito algo que analice mi portafolio específico y me sugiera acciones concretas. \\
182 & No sé. Creo que con ver todo junto ya me cambia la vida. \\
183 & Lo que más me interesa es poder medir mi rendimiento real ajustado por inflación y tipo de cambio. Hoy eso requiere cálculos manuales muy tediosos. \\
184 & Me gustaría que me digan si estoy bien o mal y por qué no solo mostrarme números sin explicación \\
185 & El onboarding de casi todas las apps de inversión es un asco, muy técnico para alguien que empieza \\
186 & Necesito poder ver mi exposición real por moneda, sector y región en un solo lugar. Hoy eso requiere mucho trabajo manual. \\
187 & Sería muy útil poder simular escenarios: qué pasa con mi portafolio si el dólar sube un 30\% o si baja el riesgo país. \\
188 & Necesito algo que me muestre todo junto en un solo lugar y que se actualice solo sin que yo tenga que hacer nada \\
189 & Consolidación automática sin cargar nada a mano. \\
190 & Valoro mucho la transparencia en las recomendaciones. Si una IA me sugiere algo, necesito ver las fuentes y el razonamiento antes de actuar. \\

\end{longtable}
}
